\subsection{Labeling Horizon and Step Size}
\label{sec:labeling-horizon}

The Receding Horizon Estimator slides a window of $N$ steps at spacing
$\mathrm{d}t$ over the synchronised commands and odometry, fits a single
parameter set $(\mu,\nu,\delta\theta)$ per window, and advances the window by
$\textsc{step\_size}$. This labeling grid is distinct from the network
prediction horizon $H$ of Section~\ref{sec:horizon}.

We set $\mathrm{d}t = 0.1\,$s to match the recorded command and odometry rate of the wayfaster dataset.

The window length $N = 10$ spans $1.0\,$s, which is the main trade-off:
it must be long enough that the robot accumulates enough translation and rotation for $(\mu,\nu)$ to be observable,
yet short enough so that the average command exitation actually represents the window.
On real missions the labeling error is flat for $N \in [8,16]$ at this $\mathrm{d}t$,
so the choice is not sensitive.

Finally $\textsc{step\_size} = 2$ emits a new
label every $0.2\,$s, matching the label rate of the original WayFASTER dataset;
the resulting window overlap of $N-2$ steps smooths consecutive estimates.
